% RÉSUMÉ
\chapter*{RÉSUMÉ}
\addcontentsline{toc}{chapter}{Résumé}
\thispagestyle{plain}
\setlength{\parskip}{0.45em}

Ce mémoire porte sur la planification des trajectoires d'un robot parallèle à câbles destiné à déplacer une caméra au-dessus d'une parcelle agricole. L'objectif est de couvrir les points de prise de vue tout en respectant les contraintes mécaniques du robot et en réduisant le temps de déplacement ainsi que l'énergie mécanique.

Une chaîne logicielle a été développée en Python. Elle génère les points de la mission, vérifie leur faisabilité dans l'espace de travail faisable en effort (WFW), calcule les coûts de transition et optimise l'ordre de visite avec l'algorithme multiobjectif NSGA-II. Les trajectoires sont ensuite profilées avec un modèle LSPB, validées numériquement et comparées à une baseline en zigzag.

Dans la zone sécurisée, les grilles $7\times7$ et $9\times9$ conservent respectivement $49/49$ et $81/81$ points. Par rapport au zigzag, l'optimisation réduit le temps de déplacement de $22{,}7\%$ et $21{,}0\%$, et l'énergie mécanique estimée de $27{,}6\%$ et $28{,}7\%$. La zone étendue $9\times9$ entraîne le rejet de quatre points et une forte augmentation de l'énergie. La zone sécurisée est donc retenue comme meilleur compromis entre couverture, faisabilité et coût de déplacement.

Ces résultats sont obtenus par simulation avec un modèle de câbles rigides et une énergie quasi-statique. Une validation sur le prototype reste nécessaire avant l'exploitation opérationnelle des trajectoires et des images par le SP3.

\vspace{0.25cm}
\textbf{Mots-clés :} robot parallèle à câbles, CDPR, planification de trajectoire, WFW, NSGA-II, optimisation multiobjectif, agriculture de précision.
